%% RESUMO NA LINGUA VERNACULA (elemento obrigatorio -- NBR 6028)
%%
%% Regras:
%%   - paragrafo unico, sem recuo, justificado;
%%   - Times New Roman 12, espacamento simples;
%%   - entre 150 e 500 palavras;
%%   - verbo na voz ativa e na terceira pessoa do singular;
%%   - as palavras-chave sao definidas em config/dados.tex.
%%
%% Um bom resumo informa: objetivo geral; principais conceitos e autores da
%% fundamentacao teorica; etapas e acoes da execucao; populacao envolvida e
%% periodo; principais resultados alcancados; conclusoes.

\begin{resumo}
A pesquisa trata da necessidade de produzir dados sintéticos de risco
cibernético que sejam suficientemente fiéis e úteis para análise, sem pressupor
que a geração sintética, por si só, assegure privacidade ou validade analítica.
O pipeline metodológico será organizado em seis etapas: análise exploratória e
preparação dos dados; geração sintética com CTGAN e, quando viável, um baseline
não neural; avaliação de fidelidade; avaliação de utilidade analítica; avaliação
de risco de exposição; e síntese dos resultados e definição das condições de
uso. A utilidade será avaliada pelo protocolo TSTR, comparado a TRTR, em uma
tarefa analítica definida a priori. O artefato será considerado tecnicamente
promissor se apresentar, simultaneamente, evidências de fidelidade, utilidade
analítica adequada ao workload definido e ausência de sinais relevantes de
exposição.
\end{resumo}
